\documentclass[12pt,a4paper]{article}

\usepackage{cmap}
\usepackage[T2A]{fontenc}
\usepackage[utf8]{inputenc}
\usepackage[english,russian]{babel}
\usepackage[a4paper,left=2.5cm,right=2cm,top=2cm,bottom=2cm]{geometry}
\usepackage{longtable}
\usepackage{booktabs}
\usepackage{array}
\usepackage{enumitem}
\usepackage{amsmath}
\usepackage[table]{xcolor}
\usepackage{hyperref}
\hypersetup{
  colorlinks=true,
  linkcolor=blue!50!black,
  urlcolor=blue!50!black,
  pdftitle={ТЗ на веб-версию игры «Конфетная фабрика»},
  pdfauthor={Техническое задание}
}

\setlength{\parindent}{0pt}
\setlength{\parskip}{6pt}

\newcommand{\field}[1]{\texttt{#1}}
\newcommand{\btn}[1]{«#1»}

\title{\textbf{Техническое задание}\\[4pt]
Веб-приложение «Конфетная фабрика»\\[4pt]
\large (онлайн-монитор командного турнира по решению задач\\
с автоматическим учётом решений через informatics.msk.ru)}
\author{Версия 1.0}
\date{3 июля 2026 г.}

\begin{document}

\maketitle

\tableofcontents
\newpage

%=====================================================================
\section{Общие сведения}

\subsection{Что делаем}

Веб-версия игры-соревнования «Конфетная фабрика». Существующая эталонная
реализация --- настольное приложение на Lazarus (описано в отдельном ТЗ
«Игра-соревнование ``Конфетная фабрика''», файл \field{TZ-candy-factory.pdf});
все игровые правила и формулы наследуются оттуда с обобщениями, явно
описанными в разделе~\ref{sec:rules}.

Ключевые отличия веб-версии от настольной:

\begin{enumerate}[nosep]
  \item Приложение разворачивается на сервере и доступно через браузер.
  \item Есть \textbf{админка} по адресу \field{/admin} с простейшей
        авторизацией (логин/пароль в файле на диске).
  \item Из админки создаются \textbf{игры}; несколько игр могут идти
        \textbf{параллельно и независимо}.
  \item При создании игры задаются ссылки на задачи с
        \url{https://informatics.msk.ru/} или \url{https://informatics.mccme.ru/}
        (это зеркала одного сайта) и список команд с их
        идентификаторами пользователей (аккаунтов) на информатиксе.
  \item У каждой команды --- \textbf{своя страница} (за логином/паролем,
        заданными в конфигурации игры) с таблицей её задач, где ячейки ---
        ссылки на условия задач. Внутри каждой строки таблицы (строка =
        уровень сложности) задачи \textbf{перемешаны} индивидуально для
        каждой команды.
  \item Есть \textbf{публичная страница игры} --- аналог старого табло
        (запасы, скорость, сетка задач с цветами состояний, гистограмма,
        таймер), но \textbf{без ссылок} на задачи. Доступна всем без
        авторизации.
  \item Решения задач подтягиваются \textbf{автоматически} с информатикса
        по аккаунтам команд (механизм заимствуется из проекта
        \url{https://github.com/FynjyBath/standings-edu} и описан в
        разделе~\ref{sec:informatics} настолько подробно, что обращаться к
        репозиторию не обязательно).
  \item Все действия в игре (покупки задач, покупки тестов, решения)
        хранятся как \textbf{журнал событий} \((\text{действие},
        \text{время})\). Админ может добавлять, редактировать и удалять
        события (кроме автоматических решений --- их можно только отключать),
        и всё состояние игры детерминированно пересчитывается из журнала.
        Страницы участников и зрителей подхватывают изменения автоматически.
\end{enumerate}

\subsection{Роли}

\begin{longtable}{@{}p{3.2cm}p{12.3cm}@{}}
\toprule
\textbf{Роль} & \textbf{Возможности} \\
\midrule
Администратор (ведущий) & Вход в \field{/admin}. Создание/настройка/запуск/завершение игр, ведение журнала событий (покупки и т.\,п.), просмотр аномалий, просмотр любых страниц. \\
Команда & Вход на свою страницу игры по логину/паролю команды. Видит общее табло и свою таблицу задач со ссылками на условия и статусами. Никаких действий не совершает --- только смотрит и решает задачи на информатиксе. \\
Зритель & Публичная страница игры без авторизации: табло без ссылок на задачи. \\
\bottomrule
\end{longtable}

Покупки задач и тестов, как и в настольной версии, вводит администратор
(команды сообщают о желании купить голосом/чатом). Решения задач
администратор не вводит --- они приходят автоматически с информатикса.

%=====================================================================
\section{Архитектура и технологический стек}
\label{sec:arch}

\subsection{Состав}

Одно серверное приложение (монолит):

\begin{itemize}[nosep]
  \item HTTP-сервер: HTML-страницы (серверный рендеринг шаблонами) +
        JSON-API для автообновления;
  \item фоновый \textbf{опросчик} (poller) информатикса --- одна горутина/поток
        на процесс, обходит аккаунты всех активных игр;
  \item хранилище: \textbf{SQLite} (один файл БД) + файлы конфигурации и
        кеша в каталоге данных.
\end{itemize}

\subsection{Стек (обязательный)}

\begin{itemize}[nosep]
  \item Язык --- \textbf{Go} (1.22+). Причина: клиент информатикса
        переносится из standings-edu (Go) почти без изменений.
  \item HTTP --- стандартный \field{net/http}, шаблоны ---
        \field{html/template}.
  \item SQLite --- драйвер без cgo: \field{modernc.org/sqlite}.
  \item Фронтенд --- обычный HTML/CSS + ванильный JavaScript
        (периодический \field{fetch} JSON-состояния). Без сборщиков,
        без SPA-фреймворков.
  \item Итог сборки --- один статический бинарник; статика и шаблоны
        встраиваются через \field{embed}.
\end{itemize}

\subsection{Каталог данных}

Путь задаётся флагом/переменной окружения (\field{--data-dir},
по умолчанию \field{./data}):

\begin{small}\begin{verbatim}
data/
  app.db                        # SQLite
  credentials/
    admin_credentials.json      # {"login": "...", "password": "..."}
    informatics_credentials.json# сервисный аккаунт информатикса
  cache/
    informatics_state.json      # кеш опросчика (см. 7.6)
  server.log                    # текстовый лог приложения
\end{verbatim}\end{small}

Файл \field{admin\_credentials.json} --- «самый простой вход»: логин и пароль
открытым текстом, читается при каждой попытке входа (правку файла можно
делать без рестарта):

\begin{small}\begin{verbatim}
{ "login": "admin", "password": "change_me" }
\end{verbatim}\end{small}

%=====================================================================
\section{Игровые правила}
\label{sec:rules}

Правила полностью повторяют настольную версию, но обобщены с трёх уровней
сложности до $n$ уровней.

\subsection{Структура задач: $n$ уровней, $n^2$ задач}

\begin{itemize}[nosep]
  \item В конфигурации игры задаётся число уровней сложности $n$
        (целое, $2 \le n \le 6$; значение по умолчанию $n=3$).
  \item Задаётся ровно $n^2$ ссылок на задачи информатикса: $n$ строк
        (строка $r$ --- уровень сложности $r$, нумерация от простого к
        сложному) по $n$ задач в строке.
  \item Набор задач \textbf{общий для всех команд}. Но порядок задач
        \textbf{внутри каждой строки} для каждой команды свой: при запуске
        игры для каждой пары (команда, строка) генерируется и сохраняется
        случайная перестановка $n$ задач этой строки. Перестановка
        генерируется один раз и далее не меняется (в том числе при
        рестарте сервера).
  \item Ячейки таблицы команды нумеруются сквозными номерами
        $1 \dots n^2$ по строкам: ячейка $(r, c)$ имеет номер $(r-1)\cdot n + c$.
        Из-за перемешивания одна и та же задача у разных команд стоит в
        разных ячейках \emph{одной и той же строки} --- зрители и соперники
        не могут по номеру ячейки понять, какая именно это задача.
\end{itemize}

\subsection{Параметры игры}

Все параметры задаются в конфигурации игры. Для каждого уровня $r = 1..n$:

\begin{longtable}{@{}p{5.2cm}p{3.2cm}p{6.8cm}@{}}
\toprule
\textbf{Параметр} & \textbf{Обозначение} & \textbf{Значение по умолчанию ($n=3$)} \\
\midrule
Цена задачи & $\mathrm{TaskCost}_r$ & 12\,000 / 12\,000 / 12\,000 \\
Цена теста (подсказки) & $\mathrm{TestCost}_r$ & 3\,000 / 7\,000 / 10\,000 \\
Нагрузка от задачи & $\mathrm{Load}_r$ & 2 / 2 / 2 \\
Бонус к запасам & $\mathrm{AmountBonus}_r$ & 12\,000 / 25\,000 / 50\,000 \\
Бонус к скорости & $\mathrm{SpeedBonus}_r$ & 4 / 7 / 11 \\
\bottomrule
\end{longtable}

Общие параметры игры:

\begin{longtable}{@{}p{5.2cm}p{3.6cm}p{6.4cm}@{}}
\toprule
\textbf{Параметр} & \textbf{Обозначение} & \textbf{По умолчанию} \\
\midrule
Стартовые запасы & $\mathrm{StartAmount}$ & 20\,000 \\
Стартовая скорость & $\mathrm{StartSpeed}$ & 15 \\
Длительность, секунд & $\mathrm{Duration}$ & 5\,100 (1 ч 25 мин) \\
\bottomrule
\end{longtable}

\subsection{Состояния задачи и механики}

У каждой пары (команда, задача) --- состояние: \field{hidden} (не куплена,
серая), \field{bought} (куплена, жёлтая), \field{passed} (решена, зелёная)
и счётчик купленных тестов. Механики (эффекты применяются в момент времени
события, см.~раздел~\ref{sec:events}):

\begin{itemize}[nosep]
  \item \textbf{Покупка задачи} (\field{hidden} $\to$ \field{bought}):
        $\mathrm{amount} \mathrel{-}= \mathrm{TaskCost}_r$;
        $\mathrm{speed} \mathrel{-}= \mathrm{Load}_r$.
  \item \textbf{Покупка теста} (только у \field{bought}):
        $\mathrm{amount} \mathrel{-}= \mathrm{TestCost}_r$; счётчик тестов $+1$.
        Число тестов не ограничено. Состояние не меняется.
  \item \textbf{Решение задачи} (\field{bought} $\to$ \field{passed}):
        $\mathrm{amount} \mathrel{+}= \mathrm{AmountBonus}_r$;
        $\mathrm{speed} \mathrel{+}= \mathrm{SpeedBonus}_r + \mathrm{Load}_r$
        (нагрузка возвращается, бонус добавляется).
  \item Потраченное на задачу и тесты при решении \textbf{не возвращается}.
  \item \textbf{Производство}: раз в секунду игрового времени
        $\mathrm{amount} \mathrel{+}= \mathrm{speed}$ (формально --- см.
        \ref{sec:statecalc}).
  \item Побеждает команда с наибольшим $\mathrm{amount}$ на момент конца игры.
\end{itemize}

Отдельных операций «отмена» больше нет: ошибочные действия исправляются
редактированием журнала событий в админке (раздел~\ref{sec:events}).

При вводе покупки админка \emph{предупреждает}, если у команды на этот момент
недостаточно средств ($\mathrm{amount} < \mathrm{TaskCost}_r$ /
$\mathrm{TestCost}_r$) или производительности
($\mathrm{speed} < \mathrm{Load}_r$), --- тексты предупреждений
«Недостаточно средств», «Недостаточно производительности», --- но позволяет
сохранить событие после явного подтверждения (чекбокс/повторное нажатие):
администратор --- окончательный авторитет.

\subsection{Жизненный цикл игры}

\begin{center}
\begin{tabular}{@{}llll@{}}
\toprule
\textbf{Статус} & \textbf{Как попадает} & \textbf{Что можно менять} & \textbf{Опрос информатикса} \\
\midrule
\field{draft} & создание игры & всё & нет \\
\field{running} & кнопка \btn{Старт} / наступление \field{start\_at} & события; продление $\mathrm{Duration}$; пароли команд & да \\
\field{finished} & $now \ge \mathrm{start\_at} + \mathrm{Duration}$ & события (временем $\le$ конца игры) & ещё 15 минут после конца, затем нет \\
\field{archived} & кнопка \btn{В архив} & ничего & нет \\
\bottomrule
\end{tabular}
\end{center}

Кнопка \btn{Старт} устанавливает $\mathrm{start\_at} := now$; альтернативно
админ может заранее задать плановое время старта, и игра перейдёт в
\field{running} автоматически. Переход \field{running} $\to$ \field{finished}
происходит сам по времени; статус --- вычислимая функция времени, отдельного
«процесса завершения» нет. Состав задач, $n$, перестановки и параметры
уровней после старта \textbf{не редактируются}.

Опрос информатикса продолжается 15 минут после конца игры, чтобы подтянуть
посылки, отправленные «на флажке» и протестированные с задержкой; решения
с временем посылки $\le$ конца игры засчитываются, позже --- попадают в
аномалии (см.~\ref{sec:autoevents}).

%=====================================================================
\section{Модель данных}
\label{sec:data}

Схема SQLite (типы условные; \field{TEXT}-время --- RFC3339 UTC):

\begin{small}\begin{verbatim}
games(id, title, status_archived INT, n INT,
      start_amount INT, start_speed INT, duration_sec INT,
      start_at TEXT NULL, created_at TEXT)

game_levels(game_id, level INT,          -- 1..n
      task_cost INT, test_cost INT, load INT,
      amount_bonus INT, speed_bonus INT,
      PRIMARY KEY(game_id, level))

game_tasks(id, game_id, level INT, ord INT, -- ord: 1..n внутри уровня,
      chapter_id INT,                       -- как ввёл админ
      url TEXT,                             -- нормализованная ссылка
      UNIQUE(game_id, chapter_id))

teams(id, game_id, name TEXT,
      informatics_user_id INT,             -- id аккаунта на информатиксе
      login TEXT, password TEXT,           -- вход на страницу команды
      UNIQUE(game_id, login))

team_task_order(team_id, cell INT,         -- cell: 1..n^2
      task_id,                             -- -> game_tasks.id
      PRIMARY KEY(team_id, cell))          -- материализованная перестановка

events(id, game_id, team_id, task_id NULL,
      type TEXT,               -- 'buy_task' | 'buy_test' | 'solve'
      at TEXT,                 -- момент события, UTC
      source TEXT,             -- 'manual' | 'auto'
      run_id INT NULL,         -- id посылки информатикса (для auto)
      enabled INT DEFAULT 1,
      comment TEXT DEFAULT '',
      created_at TEXT, updated_at TEXT,
      UNIQUE(game_id, run_id))            -- дедупликация auto-решений

anomalies(id, game_id, team_id, task_id, run_id INT,
      run_at TEXT, reason TEXT, resolved INT DEFAULT 0,
      created_at TEXT, UNIQUE(game_id, run_id))
\end{verbatim}\end{small}

Примечания:

\begin{itemize}[nosep]
  \item Пароли команд хранятся открытым текстом (осознанное упрощение:
        ценность --- нулевая, турнир одноразовый).
  \item \field{team\_task\_order} заполняется при переходе в \field{running}:
        для каждой команды и каждого уровня $r$ --- случайная перестановка
        задач уровня, записанная в ячейки $(r-1)\,n+1 \dots r\,n$.
        Генератор --- криптографический или seeded \field{math/rand};
        главное --- результат сохраняется и стабилен.
  \item Сессии (админа и команд) можно не хранить в БД: подписанные
        cookie (HMAC-SHA256 с серверным секретом, генерируемым при первом
        запуске в \field{data/secret}) с полями (роль, game\_id, team\_id,
        срок действия 24 ч).
\end{itemize}

%=====================================================================
\section{Журнал событий и расчёт состояния}
\label{sec:events}

\subsection{Принцип: состояние = функция журнала}

Никакие «текущие запасы/скорости» в БД не хранятся. Всё состояние игры в
любой момент $t$ детерминированно вычисляется из конфигурации игры и
списка \textbf{включённых} (\field{enabled=1}) событий с $\field{at} \le t$.
Поэтому редактирование прошлого (изменение времени события, удаление,
добавление задним числом) автоматически и корректно меняет текущие цифры
на всех страницах.

\subsection{Типы событий}

\begin{longtable}{@{}p{2.6cm}p{2.4cm}p{10.5cm}@{}}
\toprule
\textbf{Тип} & \textbf{Источник} & \textbf{Смысл и эффекты} \\
\midrule
\field{buy\_task} & manual & Покупка задачи командой. $-\mathrm{TaskCost}_r$ к запасам, $-\mathrm{Load}_r$ к скорости, задача $\to$ \field{bought}. \\
\field{buy\_test} & manual & Покупка теста. $-\mathrm{TestCost}_r$ к запасам, счётчик тестов задачи $+1$. \\
\field{solve} & auto (poller) или manual & Решение задачи. $+\mathrm{AmountBonus}_r$ к запасам, $+(\mathrm{SpeedBonus}_r + \mathrm{Load}_r)$ к скорости, задача $\to$ \field{passed}. Auto-событие ссылается на \field{run\_id} посылки. \\
\bottomrule
\end{longtable}

Правила редактирования в админке:

\begin{itemize}[nosep]
  \item ручные события (\field{source=manual}): добавлять, менять время,
        команду, задачу, тип; удалять физически;
  \item автоматические (\field{source=auto}): удалять и менять
        нельзя (иначе опросчик создаст их заново); можно
        \textbf{отключить} (\field{enabled=0}) --- событие остаётся в списке
        зачёркнутым и не участвует в расчёте, а уникальность по
        \field{run\_id} гарантирует, что опросчик не создаст дубль; можно
        \textbf{поправить время} (поле \field{at}) --- связь с
        \field{run\_id} сохраняется;
  \item ручное событие \field{solve} допустимо (страховка на случай
        недоступности информатикса).
\end{itemize}

\subsection{Алгоритм расчёта состояния}
\label{sec:statecalc}

Обозначим $t_0 = \mathrm{start\_at}$,
$t_{\mathrm{end}} = t_0 + \mathrm{Duration}$,
$T = \min(now, t_{\mathrm{end}})$.

\begin{enumerate}[nosep]
  \item Берём включённые события игры с $\field{at} \le T$, сортируем по
        (\field{at}, \field{id}). События с $\field{at} < t_0$ трактуются
        как произошедшие в $t_0$ (валидатор их подсвечивает).
  \item Для каждой команды ведём $(\mathrm{amount}, \mathrm{speed}) :=
        (\mathrm{StartAmount}, \mathrm{StartSpeed})$ и состояния задач.
  \item Идём по целым секундам $k = 1, 2, \dots, \lfloor T - t_0 \rfloor$
        и по событиям в хронологическом порядке, применяя на каждом шаге
        сначала все события с $\field{at} \le t_0 + k$, затем начисление
        $\mathrm{amount} \mathrel{+}= \mathrm{speed}$ за секунду $k$.
        (Реализация обязана быть эквивалентной этому определению, но
        эффективной: между событиями скорость постоянна, поэтому суммарное
        начисление за интервал считается умножением, а не циклом.)
  \item Эффекты события применяются согласно таблице типов. Событие
        «не по правилам» (см. валидацию ниже) эффект всё равно применяет ---
        админ отвечает за журнал.
\end{enumerate}

Результат расчёта: для каждой команды --- $\mathrm{amount}$,
$\mathrm{speed}$, по каждой ячейке --- состояние и число тестов.

\textbf{Валидация} (не блокирует, а подсвечивает в админке список проблем):

\begin{itemize}[nosep]
  \item \field{buy\_task} по задаче не в \field{hidden};
        \field{buy\_test} по задаче не в \field{bought};
        \field{solve} по задаче не в \field{bought};
  \item в какой-то момент $\mathrm{amount} < 0$ или $\mathrm{speed} < 0$;
  \item событие раньше $t_0$ или позже $t_{\mathrm{end}}$.
\end{itemize}

Расчёт кешируется в памяти процесса: пересчитывать при (а) изменении
журнала/конфигурации, (б) не реже раза в секунду для начислений (достаточно
пересчитывать лениво при запросе состояния, если с прошлого расчёта прошло
$\ge 1$ с или изменились данные).

%=====================================================================
\section{Ссылки на задачи информатикса}
\label{sec:urls}

Админ вводит $n^2$ ссылок вида:

\begin{small}\begin{verbatim}
https://informatics.msk.ru/mod/statements/view.php?chapterid=113449
https://informatics.mccme.ru/mod/statements/view.php?chapterid=2958#1
\end{verbatim}\end{small}

Требования к разбору (нормализации):

\begin{enumerate}[nosep]
  \item Допустимые хосты: \field{informatics.msk.ru},
        \field{www.informatics.msk.ru}, \field{informatics.mccme.ru},
        \field{www.informatics.mccme.ru}. Это зеркала: считаются
        \textbf{эквивалентными}.
  \item Путь: \field{/mod/statements/view.php} (без учёта регистра),
        параметр запроса \field{chapterid} --- положительное целое.
        Фрагмент (\field{\#1}) и прочие параметры игнорируются.
  \item Идентификатором задачи считается \field{chapterid} --- он же
        совпадает с \field{problem.id} в API посылок (см.~\ref{sec:runsapi}).
  \item Канонический URL для показа на страницах команд:
        \field{https://informatics.msk.ru/mod/statements/view.php?chapterid=<id>}
        (базовый хост берётся из \field{base\_url} в
        \field{informatics\_credentials.json}, чтобы можно было переключить
        зеркало одним местом).
  \item Дубликаты \field{chapterid} внутри игры запрещены --- ошибка формы.
  \item Ссылка, не подошедшая под правила, --- ошибка формы с указанием
        конкретной строки.
\end{enumerate}

Идентификатор аккаунта команды --- числовой \field{user\_id} на
информатиксе (виден в URL профиля:
\field{/user/profile.php?id=<user\_id>}). Валидация: положительное целое.

%=====================================================================
\section{Интеграция с informatics (перенос из standings-edu)}
\label{sec:informatics}

Этот раздел --- переработанное под нашу задачу описание клиента информатикса
из \url{https://github.com/FynjyBath/standings-edu}
(файл \field{internal/source/informatics.go}). Реализовать нужно именно это
поведение; код из репозитория можно брать за основу.

\subsection{Сервисный аккаунт}

Файл \field{data/credentials/informatics\_credentials.json}:

\begin{small}\begin{verbatim}
{
  "username": "...",
  "password": "...",
  "base_url": "https://informatics.msk.ru"
}
\end{verbatim}\end{small}

\field{base\_url} опционален (по умолчанию \field{https://informatics.msk.ru});
поля обрезаются от пробелов; пустые логин/пароль --- фатальная ошибка старта
опросчика (но не веб-сервера: страницы должны работать, в админке ---
заметный баннер «опрос информатикса не работает: …»).

\subsection{Авторизация (Moodle-логин)}
\label{sec:login}

HTTP-клиент: таймаут 20 с, cookie jar (сессия живёт в куках).

\begin{enumerate}[nosep]
  \item \field{GET \{base\_url\}/login/index.php}; из HTML извлечь токен
        регулярным выражением
        \verb|name="logintoken"\s+value="([^"]+)"|.
        Токен не найден --- ошибка «informatics login token not found».
  \item \field{POST} на тот же URL,
        \field{Content-Type: application/x-www-form-urlencoded},
        заголовок \field{Referer} = URL логина; поля формы:
        \field{anchor}="", \field{logintoken}, \field{username},
        \field{password}, \field{rememberusername}=0.
  \item Разбор ответа: если тело содержит \field{logout.php} --- успех.
        Если содержит \verb|name="logintoken"| или форму логина или
        \field{loginerrors} --- ошибка «bad credentials or blocked account».
\end{enumerate}

Повторный логин: сессия считается живой 30 минут с последнего успешного
логина; логин выполняется заново (а) по истечении этого срока, (б) немедленно
при ответе API «Not authorized» (см. ниже), после чего запрос повторяется
один раз. Логин защищён мьютексом --- параллельные запросы не должны
логиниться наперегонки.

\subsection{API посылок}
\label{sec:runsapi}

Список посылок пользователя, новые сверху, постранично:

\begin{small}\begin{verbatim}
GET {base_url}/py/problem/0/filter-runs
    ?problem_id=0&user_id={user_id}&count=1000&page={page}
    &from_timestamp=-1&to_timestamp=-1&lang_id=-1
    &status_id=-1&statement_id=0&with_comment=
\end{verbatim}\end{small}

Ответ --- JSON:

\begin{small}\begin{verbatim}
{
  "result": "success" | "error",
  "message": "...",
  "data": [ {
      "id": 123456,              // id посылки, монотонно растет
      "create_time": "...",      // время посылки
      "ejudge_status": 0,        // вердикт ejudge
      "ejudge_score": 100,       // может быть null/строкой/числом
      "score": 100,              // то же
      "problem": { "id": 113449 } // == chapterid задачи
  }, ... ],
  "metadata": { "page_count": 3, "count": 2456 }
}
\end{verbatim}\end{small}

Обязательные детали обработки:

\begin{itemize}[nosep]
  \item HTTP-статус $\ne$ 200 --- ошибка (в лог, с телом до 1 КБ).
  \item \field{result="error"} и \field{message="Not authorized"}
        (без учёта регистра) --- перелогин и один повтор
        (\ref{sec:login}); другой \field{error} --- ошибка с текстом
        \field{message}.
  \item \field{ejudge\_score}/\field{score} в живых данных бывают числом,
        строкой с числом, пустой строкой и \field{null} --- парсер обязан
        принимать все варианты (нечисловое $\Rightarrow$ «нет значения»).
  \item \textbf{Решённой} посылка считается при
        \field{ejudge\_status} $\in \{0, 8\}$: $0$ = OK, $8$ = ACCEPTED
        («Зачтено»). Остальные вердикты для нас --- просто попытки, они
        нигде не учитываются.
  \item Поле \field{create\_time}: standings-edu его не разбирает, нам оно
        \textbf{нужно} (время события \field{solve}). Исполнителю
        необходимо посмотреть на реальный ответ API и определить точный
        формат (ожидается строка даты-времени по московскому времени или
        unix-время). Распарсенное время приводится к UTC. Если поле
        отсутствует/не парсится --- взять момент обнаружения посылки
        опросчиком и записать предупреждение в лог.
  \item Размер страницы 1000; страницы, начиная со второй, можно качать
        параллельно (в standings-edu --- до 8 воркеров), но для наших
        объёмов (свежие аккаунты команд) достаточно последовательной
        загрузки --- параллелизм опционален.
\end{itemize}

\subsection{Инкрементальная выгрузка}

Для каждого аккаунта (\field{user\_id}) хранится \field{max\_run\_id} ---
наибольший виденный \field{id} посылки. Алгоритм
\field{FetchNewRuns(user\_id)}:

\begin{enumerate}[nosep]
  \item Загрузить страницу 1. Идя по посылкам (они отсортированы по убыванию
        \field{id}), собирать посылки с $\field{id} > \field{max\_run\_id}$;
        встретив $\field{id} \le \field{max\_run\_id}$ --- остановиться
        (дальше всё старое).
  \item Если не встретили старых и \field{page\_count} > 1 --- продолжить
        со страницы 2 и далее до первой старой посылки или конца.
  \item Если \field{max\_run\_id} ещё не было (первый опрос аккаунта) ---
        выкачать все страницы.
  \item Обновить \field{max\_run\_id} и сохранить кеш (см.~\ref{sec:cache}).
  \item Вернуть новые посылки.
\end{enumerate}

\subsection{Опросчик}

Один фоновый цикл с периодом \field{poll\_interval} (конфиг, по умолчанию
\textbf{30 с}):

\begin{enumerate}[nosep]
  \item Составить множество аккаунтов: все команды всех игр в статусе
        \field{running}, плюс \field{finished} менее 15 минут назад.
        Пусто --- спать до следующего цикла.
  \item Для каждого аккаунта последовательно (пауза $\ge 1$ с между
        аккаунтами, чтобы не долбить сайт): \field{FetchNewRuns}.
  \item Каждую новую \textbf{решённую} посылку сматчить: по
        \field{problem.id} найти задачу игры (\field{game\_tasks.chapter\_id}),
        по \field{user\_id} --- команду. Посылки по чужим задачам ---
        игнорировать молча.
  \item Для сматченной посылки со временем $t_{run}$ (из \field{create\_time}):
        \begin{itemize}[nosep]
          \item если по паре (команда, задача) уже есть включённое событие
                \field{solve} или запись с тем же \field{run\_id} (событие
                или аномалия) --- пропустить;
          \item если на момент $t_{run}$ задача в состоянии \field{bought}
                и $t_0 \le t_{run} \le t_{\mathrm{end}}$ --- создать событие
                \field{solve} (\field{source=auto}, \field{at}=$t_{run}$,
                \field{run\_id});
          \item иначе (не куплена ещё/уже решена/вне времени игры) ---
                создать запись в \field{anomalies} с причиной; админ решит
                вручную (см.~\ref{sec:autoevents}).
        \end{itemize}
  \item Ошибка сети/логина на любом шаге: записать в лог, показать баннер в
        админке (время и текст последней ошибки), \textbf{не} терять
        накопленный кеш, повторить в следующем цикле. Один упавший аккаунт
        не должен прерывать обход остальных.
\end{enumerate}

\subsection{Кеш опросчика}
\label{sec:cache}

Файл \field{data/cache/informatics\_state.json} (схема --- упрощение
standings-edu):

\begin{small}\begin{verbatim}
{
  "version": 1,
  "accounts": {
    "285757": { "max_run_id": 1234567,
                "updated_at": "2026-07-03T12:00:00Z" }
  }
}
\end{verbatim}\end{small}

\begin{itemize}[nosep]
  \item Запись --- \textbf{атомарная}: сериализовать во временный файл в том
        же каталоге, затем \field{rename} поверх целевого (как в
        standings-edu).
  \item При чтении: файла нет --- пустой кеш; не парсится или
        \field{version} не совпадает --- игнорировать и начать с нуля
        (это лишь ускорение, источник истины --- информатикс и БД).
  \item Доступ под мьютексом.
  \item Внимание: сброс кеша приводит к повторной выкачке всех посылок, но
        \textbf{не} к дублям событий --- дедупликация по \field{run\_id}
        в БД (уникальный индекс).
\end{itemize}

\subsection{Аномалии}
\label{sec:autoevents}

Страница игры в админке показывает нерешённые аномалии --- решённые посылки,
которые нельзя было применить автоматически:

\begin{itemize}[nosep]
  \item \field{not\_bought} --- на момент посылки задача не куплена
        (например, админ ещё не внёс покупку);
  \item \field{already\_passed} --- задача уже решена;
  \item \field{out\_of\_time} --- посылка вне интервала игры.
\end{itemize}

По каждой аномалии доступны действия: \btn{Принять} (создаёт ручное событие
\field{solve} с временем посылки, аномалия помечается решённой; полезно
после того, как админ внёс пропущенную покупку) и \btn{Отклонить}
(помечается решённой без событий). Пока аномалия не решена, опросчик не
создаёт по этой посылке ничего нового (уникальность по \field{run\_id}).

%=====================================================================
\section{Страницы и интерфейс}
\label{sec:pages}

Все страницы --- на русском, вёрстка простая (таблицы/флексы), без
брендинга; должна нормально выглядеть на проекторе (публичная страница) и на
ноутбуке/телефоне (страница команды).

\subsection{Карта маршрутов}

\begin{longtable}{@{}p{6.4cm}p{2.2cm}p{6.6cm}@{}}
\toprule
\textbf{Маршрут} & \textbf{Доступ} & \textbf{Содержимое} \\
\midrule
\field{GET /} & все & Список неархивных игр со ссылками на публичные страницы. \\
\field{GET /g/\{gameId\}} & все & Публичное табло игры. \\
\field{GET /g/\{gameId\}/team} & форма входа & Логин команды; после входа --- редирект на страницу команды. \\
\field{GET /g/\{gameId\}/team/\{teamId\}} & команда или админ & Страница команды. \\
\field{GET /api/g/\{gameId\}/state} & все & JSON состояния (без ссылок). \\
\field{GET /api/g/\{gameId\}/team/\{teamId\}/state} & команда или админ & JSON состояния + ссылки и порядок задач команды. \\
\field{GET,POST /admin} и \field{/admin/...} & админ & Админка (см.~\ref{sec:admin}). \\
\bottomrule
\end{longtable}

\field{gameId}, \field{teamId} --- целые id из БД. Неизвестный id --- 404.
Страница команды под чужой командной сессией --- 403.

\subsection{Публичное табло \field{/g/\{gameId\}}}

Воспроизводит настольный монитор:

\begin{enumerate}[nosep]
  \item Заголовок: название игры; крупный таймер \btn{Оставшееся время:
        ЧЧ:ММ:СС} (до старта --- \btn{До начала: ЧЧ:ММ:СС} или
        \btn{Игра не началась}; после конца --- \btn{Игра завершена}).
  \item По колонке на команду: имя; \btn{Запасы: N}; \btn{Скорость: N};
        сетка $n \times n$ ячеек с номерами $1..n^2$, раскрашенных по
        состоянию (серый/жёлтый/зелёный; на жёлтой ячейке --- бейдж с числом
        купленных тестов, если $> 0$). \textbf{Без ссылок и без названий
        задач.}
  \item Гистограмма запасов: горизонтальные или вертикальные полосы,
        длина пропорциональна запасам; масштаб --- по максимуму среди команд
        (лидер = 100\%).
  \item Индикатор скорости у каждой команды: «бегущая» полоска
        (CSS-анимация, период обратно пропорционален скорости) --- чисто
        декоративный, как в оригинале.
  \item Таблица параметров игры: строки --- уровни, столбцы --- «Цена
        задачи», «Цена подсказки», «Нагрузка от задачи», «Бонус запасы»,
        «Бонус скорость».
  \item Легенда цветов: «Не куплена» / «Куплена» / «Решена».
\end{enumerate}

\subsection{Страница команды}

Вход: форма (логин, пароль команды), проверка по \field{teams} данной игры,
cookie на 24 ч. Содержимое --- всё то же публичное табло \textbf{плюс}
собственная таблица задач $n \times n$:

\begin{itemize}[nosep]
  \item ячейка --- номер + \textbf{ссылка} на условие задачи на информатиксе
        (канонический URL, \field{target="\_blank"}), порядок ячеек --- по
        \field{team\_task\_order} этой команды;
  \item цвет фона ячейки --- состояние задачи; бейдж числа купленных тестов;
  \item подписи строк: «Уровень 1» … «Уровень $n$» (при $n=3$ можно
        Easy/Middle/Hard);
  \item ссылки видны и до покупки задачи (состояние \field{hidden}) ---
        команда должна видеть условия, чтобы выбирать, что покупать.
\end{itemize}

Чужие таблицы задач (и чужой порядок) команде не показываются никогда.

\subsection{Автообновление}
\label{sec:refresh}

\begin{itemize}[nosep]
  \item Страницы табло и команды раз в \textbf{3 секунды} (конфиг)
        запрашивают свой \field{/api/.../state} и перерисовывают числа,
        цвета ячеек, гистограмму без перезагрузки страницы.
  \item Таймер тикает на клиенте каждую секунду; ответ API содержит
        \field{server\_time}, по нему клиент корректирует смещение своих
        часов.
  \item Ошибка запроса --- молча повторить в следующий цикл; после
        3 подряд ошибок показать ненавязчивый индикатор «нет связи».
  \item Обычная перезагрузка страницы (F5) в любом случае даёт актуальное
        состояние --- состояние всегда считается на сервере из журнала.
\end{itemize}

Формат \field{GET /api/g/\{id\}/state} (командный вариант дополнительно
содержит \field{url} в ячейках и порядок своей таблицы):

\begin{small}\begin{verbatim}
{
  "server_time": "2026-07-03T12:00:00Z",
  "status": "running",             // draft|running|finished|archived
  "start_at": "...", "end_at": "...",
  "n": 3,
  "levels": [ {"task_cost":12000, "test_cost":3000, "load":2,
               "amount_bonus":12000, "speed_bonus":4}, ... ],
  "teams": [ {
      "id": 1, "name": "Команда №1",
      "amount": 31337, "speed": 19,
      "cells": [ {"cell":1, "state":"bought", "tests":2}, ... ]
  }, ... ]
}
\end{verbatim}\end{small}

\subsection{Админка \field{/admin}}
\label{sec:admin}

Вход: форма логин/пароль, сверка с \field{admin\_credentials.json},
cookie-сессия 24 ч. Все прочие \field{/admin/*} без сессии --- редирект на
форму входа.

\paragraph{Список игр.} Таблица: название, статус, $n$, число команд,
старт, оставшееся время; кнопки \btn{Открыть}, \btn{В архив}. Кнопка
\btn{Создать игру}.

\paragraph{Создание/редактирование игры (\field{draft}).} Форма:

\begin{itemize}[nosep]
  \item название; $n$; параметры уровней (таблица $n$ строк, предзаполнена
        дефолтами из~\ref{sec:rules}); StartAmount, StartSpeed, Duration;
  \item задачи: $n$ полей-«текстариев» (по уровню), в каждый вставляется
        $n$ ссылок по одной в строке; при сохранении --- нормализация и
        валидация по разделу~\ref{sec:urls} с точными сообщениями об
        ошибках;
  \item команды (динамический список, минимум 2): имя,
        \field{informatics\_user\_id}, логин, пароль страницы команды;
        логины уникальны в игре;
  \item плановое время старта (опционально) и кнопка \btn{Старт сейчас}.
\end{itemize}

При старте: фиксируются перестановки, игра исчезает из редактирования
конфигурации.

\paragraph{Страница игры (основной рабочий экран ведущего).}

\begin{enumerate}[nosep]
  \item Шапка: статус, таймер, кнопки \btn{+15 минут} (продление
        Duration), \btn{В архив}; баннер ошибок опросчика (последняя ошибка
        и когда).
  \item \textbf{Живое табло} --- как публичное, но каждая ячейка кликабельна:
        клик открывает меню ячейки с кнопками \btn{Купить задачу} /
        \btn{Купить тест} (активность кнопок по текущему состоянию ячейки,
        как в настольной версии). Кнопка создаёт ручное событие с
        $\field{at}=now$; при нехватке средств/скорости --- предупреждение
        с подтверждением (см.~\ref{sec:rules}). Здесь же видно, какая
        informatics-задача стоит в ячейке (подпись с \field{chapterid} и
        ссылкой) --- админу можно.
  \item \textbf{Журнал событий}: таблица всех событий игры (время,
        команда, ячейка/задача, тип, источник, run\_id, вкл/выкл,
        комментарий), сортировка по времени, фильтры по команде и типу.
        Действия по строке: \btn{Изменить} (форма: время с точностью до
        секунды, команда, задача, тип --- для manual; только время --- для
        auto), \btn{Удалить} (manual), \btn{Отключить}/\btn{Включить}
        (auto). Кнопка \btn{Добавить событие} --- та же форма с ручным
        временем. После любого изменения --- мгновенный пересчёт и
        перерисовка табло.
  \item Панель \textbf{предупреждений валидации} (раздел~\ref{sec:statecalc}):
        список проблемных событий со ссылками на строки журнала.
  \item Секция \textbf{аномалий} (раздел~\ref{sec:autoevents}).
  \item Секция \textbf{команды}: логины/пароли страниц команд (с прямыми
        ссылками), возможность сменить пароль после старта.
\end{enumerate}

Мутации админки --- обычные POST-формы или fetch-запросы; CSRF-токен в
формах (генерируется в сессии). Одновременная работа двух админов не
поддерживается специально, но не должна ломать данные (все мутации ---
атомарные транзакции SQLite).

%=====================================================================
\section{Конфигурация и развёртывание}

\subsection{Параметры запуска}

Флаги/переменные окружения:

\begin{longtable}{@{}p{4.6cm}p{3.4cm}p{7cm}@{}}
\toprule
\textbf{Параметр} & \textbf{По умолчанию} & \textbf{Смысл} \\
\midrule
\field{--listen} & \field{:8080} & адрес HTTP-сервера \\
\field{--data-dir} & \field{./data} & каталог данных \\
\field{--poll-interval} & \field{30s} & период опроса информатикса \\
\field{--page-refresh} & \field{3s} & период автообновления страниц (отдаётся клиенту) \\
\bottomrule
\end{longtable}

\subsection{Развёртывание}

\begin{itemize}[nosep]
  \item Поставка: один бинарник (linux/amd64) + пример
        \field{data/}-каталога + unit-файл systemd (в комплекте,
        \field{Restart=always}).
  \item HTTPS --- забота reverse-proxy (nginx/caddy) перед приложением;
        приложение слушает обычный HTTP. В README --- пример конфига nginx.
  \item Все времена хранятся и передаются в UTC; отображение на страницах ---
        в часовом поясе сервера (по умолчанию Europe/Moscow --- указать в
        README про переменную \field{TZ}).
  \item При рестарте сервера состояние полностью восстанавливается из БД и
        кеша; идущие игры продолжаются без потерь (время старта в БД).
\end{itemize}

\subsection{Логирование}

Текстовый лог (\field{data/server.log} + stdout): старты/завершения игр,
каждое созданное/изменённое/удалённое событие (кто: admin/poller), каждая
аномалия, ошибки опросчика, входы в админку и на страницы команд (успех и
провал). Формат строки: время, уровень, сообщение.

%=====================================================================
\section{Нефункциональные требования}

\begin{enumerate}[nosep]
  \item Нагрузка: до 10 параллельных игр, до 8 команд в игре, до ~200
        одновременных зрителей. С учётом обновления раз в 3 с ---
        ~70 RPS лёгких JSON-запросов; ответ из кеша состояния, без
        пересчёта на каждый запрос.
  \item Целостность: журнал событий --- единственный источник истины;
        никакие операции не должны приводить к его частичной записи
        (транзакции).
  \item Опросчик не должен положить информатикс: не более одного
        запроса в секунду суммарно, инкрементальная выгрузка обязательна.
  \item Секреты (пароль сервисного аккаунта информатикса) не должны
        попадать в логи и в HTML-страницы.
  \item Код --- с юнит-тестами как минимум на: нормализацию ссылок
        (\ref{sec:urls}), расчёт состояния из журнала
        (\ref{sec:statecalc}, включая события задним числом и отключённые),
        инкрементальную логику \field{max\_run\_id}, матчинг посылок и
        генерацию аномалий.
\end{enumerate}

%=====================================================================
\section{Критерии приёмки}

\begin{enumerate}
  \item \textbf{Развёртывание.} На чистом сервере по README: бинарник +
        каталог данных + systemd; админка открывается, вход по паре из
        \field{admin\_credentials.json}; неправильный пароль --- отказ.
  \item \textbf{Создание игры.} Игра с $n=3$, 9 ссылками (в т.\,ч. часть
        с \field{informatics.mccme.ru} и с \field{\#1} на конце ---
        принимаются и нормализуются), 3 командами. Ошибочная ссылка и
        дубль задачи --- внятные ошибки формы.
  \item \textbf{Перемешивание.} После старта таблицы двух команд содержат
        одинаковые множества задач в каждой строке, но в разном порядке
        (проверить несколько игр); порядок стабилен после рестарта сервера.
  \item \textbf{Табло.} Публичная страница: запасы растут на
        StartSpeed в секунду, таймер убывает, ссылок на задачи нет в HTML
        вообще (в том числе в JSON публичного API). Страница команды ---
        те же цифры плюс кликабельные ссылки на условия.
  \item \textbf{Покупка.} Через админ-табло куплена задача уровня 1:
        у команды $-12\,000$ запасов, $-2$ скорости, ячейка жёлтая на всех
        страницах в течение 3--5 секунд без перезагрузки. Покупка при
        нехватке средств --- предупреждение, сохранение только после
        подтверждения.
  \item \textbf{Автозачёт.} Команда сдаёт задачу на информатиксе (вердикт
        OK или «Зачтено»): в течение периода опроса появляется событие
        \field{solve} с временем посылки; запасы и скорость меняются по
        правилам; ячейка зелёная. Посылка по некупленной задаче --- попадает
        в аномалии, после \btn{Принять} (и внесения покупки задним числом)
        зачитывается.
  \item \textbf{Редактирование истории.} Смещение времени покупки на
        10 минут назад увеличивает текущие запасы ровно на
        $10 \cdot 60 \cdot \mathrm{Load}_r$ (недоначисленная скорость) ---
        сверить с ручным расчётом. Удалённое событие исчезает со всех
        страниц после очередного автообновления. Отключённое auto-решение
        не возвращается опросчиком.
  \item \textbf{Параллельность.} Две одновременно идущие игры не влияют
        друг на друга (в т.\,ч. одна и та же задача информатикса в двух
        играх и разные команды одного «владельца»).
  \item \textbf{Отказоустойчивость.} При недоступном информатиксе табло и
        админка полностью работоспособны, в админке --- баннер с ошибкой;
        после восстановления связи пропущенные решения подтягиваются с
        корректным временем посылки. После \field{kill -9} и рестарта ---
        игра продолжается, дублей решений нет.
  \item \textbf{Финиш.} По истечении времени начисление останавливается,
        статус \field{finished}; посылка, отправленная за минуту до конца,
        но протестированная после, --- зачтена автоматически (окно
        15 минут).
\end{enumerate}

\end{document}
